% --- Perfil ---
\section{Perfil}
\begin{onecolentry}
	Soy graduado en Ingeniería de Telecomunicación, habiendo cursado el grado y el máster en la Universidad de Alcalá. Dentro del campo de la ingeniería, estoy especialmente interesado en la electrónica aplicada, la instrumentación y el desarrollo de hardware, habiendo orientado mi trayectoria profesional hacia el diseño, la implementación y la validación de sistemas electrónicos en entornos técnicos exigentes. He trabajado con instrumentación, herramientas CAD y plataformas embebidas, así como con tecnologías vinculadas al desarrollo de prototipos y a la integración de soluciones hardware y software, poniendo el foco en la fiabilidad experimental, la capacidad de medida y la adaptación de los sistemas a requisitos de alto rendimiento.
\end{onecolentry}

\vspace{-5pt}

% --- Educación ---
\section{Educación}

\begin{onecolentry}
	\begin{tabularx}{\textwidth}{@{}>{\raggedright\arraybackslash}X>{\raggedright\arraybackslash}p{3.6cm}>{\raggedleft\arraybackslash}p{2.4cm}@{}}
		\textbf{Máster en Ingeniería en Tecnologías de Telecomunicación} & Universidad de Alcalá & \textit{Febrero 2026}
	\end{tabularx}
	
	\vspace{0.05 cm}
	\hspace{0.2cm}Premio extraordinario
\end{onecolentry}

\vspace{0.2 cm}

\begin{onecolentry}
	\begin{tabularx}{\textwidth}{@{}>{\raggedright\arraybackslash}X>{\raggedright\arraybackslash}p{3.6cm}>{\raggedleft\arraybackslash}p{2.4cm}@{}}
		\textbf{Grado en Ingeniería en Tecnologías de Telecomunicación} & Universidad de Alcalá & \textit{Julio 2024}
	\end{tabularx}
	
	\vspace{0.05 cm}
	\hspace{0.2cm}2.º puesto de promoción y premio extraordinario
\end{onecolentry}

\vspace{-7pt}

% --- Experiencia ---
\section{Experiencia}

\begin{twocolentry}{\textit{Septiembre 2024 – Enero 2026}}
	\textbf{Plataformas técnicas e integración} \\
	Universidad de Alcalá, Alcalá de Henares
\end{twocolentry}
\begin{onecolentry}
	\begin{highlights}
		\item Diseñé y desplegué entornos técnicos reproducibles sobre contenedores y orquestadores para pruebas e integración de servicios.
		\item Desarrollé herramientas y pipelines en Python, MATLAB y C para adquisición, procesado y validación de datos en cargas intensivas.
		\item Integré flujos automatizados para despliegue y validación en entornos de experimentación y producción.
	\end{highlights}
\end{onecolentry}

\vspace{0.2 cm}

\begin{twocolentry}{\textit{Febrero 2024 – Agosto 2024}}
	\textbf{Sistemas y soporte técnico} \\
	Correos y Telégrafos S.A., Madrid
\end{twocolentry}
\begin{onecolentry}
	\begin{highlights}
		\item Participé en tareas de soporte técnico y documentación de infraestructuras y servicios corporativos.
		\item Administré sistemas Linux empresariales y colaboré en despliegues automatizados y prácticas DevOps.
		\item Gestioné servicios de seguridad y autenticación empleados en entornos empresariales.
	\end{highlights}
\end{onecolentry}

\vspace{0.2 cm}

\begin{twocolentry}{\textit{Enero 2023 – Enero 2024}}
	\textbf{Ingeniería e investigación aplicada} \\
	Universidad de Alcalá / Correos y Telégrafos S.A., Alcalá de Henares
\end{twocolentry}
\begin{onecolentry}
	\begin{highlights}
		\item Desarrollé plataformas técnicas de análisis de datos para proyectos de investigación.
		\item Implementé soluciones de computación multihilo y en clúster para acelerar procesos intensivos.
		\item Me encargué de la administración de entornos Linux y colaboré en tareas de soporte técnico y validación experimental.
		\item Participé como coautor en publicaciones científicas \href{https://orcid.org/0009-0000-6862-1872}{ORCID}.
	\end{highlights}
\end{onecolentry}

\vspace{-7pt}

% --- Aptitudes ---
\section{Aptitudes}

\begin{onecolentry}
	\textbf{Programación:} C/C++, Python, MATLAB, VHDL, Bash
\end{onecolentry}
\begin{onecolentry}
	\textbf{Instrumentación y CAD:} LabVIEW, AutoCAD, Autodesk Inventor, SketchUp, PSpice, KiCad
\end{onecolentry}
\begin{onecolentry}
	\textbf{MCUs:} Xilinx Zynq, STM32, ESP32, NXP 17xx, Renesas Synergy
\end{onecolentry}
\begin{onecolentry}
	\textbf{Electrónica y diseño:} Electrónica analógica y digital, instrumentación, DAQs de alta velocidad, diseño y validación de hardware
\end{onecolentry}